\documentclass[11pt]{article}
\usepackage[margin=1in]{geometry}

% Core packages
\usepackage{amsmath,amssymb}
\usepackage{amsthm}
\usepackage{booktabs}
\usepackage{multirow}
\usepackage{adjustbox}
\usepackage{algorithm}
\usepackage{algpseudocode}
\usepackage{natbib}
\usepackage{xcolor}
\usepackage{soul}
\usepackage{pifont}
\newcommand{\change}[1]{{\color{red}#1}}
\newcommand{\cmark}{\ding{51}}
\newcommand{\xmark}{\ding{55}}

% Theorem-like environments
\theoremstyle{plain}
\newtheorem{theorem}{Theorem}[section]
\newtheorem{lemma}[theorem]{Lemma}
\newtheorem{proposition}[theorem]{Proposition}
\newtheorem{corollary}[theorem]{Corollary}
\newtheorem{claim}[theorem]{Claim}

\theoremstyle{definition}
\newtheorem{definition}[theorem]{Definition}
\newtheorem{assumption}[theorem]{Assumption}
\newtheorem{example}[theorem]{Example}

\theoremstyle{remark}
\newtheorem{remark}[theorem]{Remark}

\setlength{\parindent}{0pt}
\setlength{\parskip}{1\baselineskip}

\title{Untitled}
\author{Author Name}
\date{\today}

\begin{document}
\maketitle

% ============================================================
%  ABSTRACT
% ============================================================
\begin{abstract}
  Many classical combinatorial algorithms follow a natural residual loop: score
  elements, prune or select one, and update the problem state. When the maintained
  statistics are additive---such as degree or marginal gain---each update costs $O(n)$,
  yielding an overall $O(n^2)$ decoder. GNNs can learn richer scoring rules, but their
  nonlinearity breaks this efficiency: small approximation errors accumulate across
  pruning steps, causing incremental state updates to diverge, and forcing a full
  network rerun after every decision---inflating complexity to $O(n^3)$.

  We ask: \textit{can a GNN be trained to behave additively, recovering $O(n^2)$
  decoding without sacrificing expressiveness?} We prove that standard GNNs cannot:
  universality implies irreducible approximation error that propagates through the
  residual loop. We then show that for any locally decomposable objective, the objective
  gradient provides a local additive correction after each vertex deletion, and that
  under non-saturating activations this signal survives nonlinear processing. This
  insight yields two complementary routes to $O(n^2)$ decoding. The first feeds the
  gradient to a lightweight learned updater, which can represent the exact residual
  correction in its first linear layer. The second injects the objective gradient into
  the nonlinear backbone and trains a final additive coordinate to track the negative
  gradient, enabling a fully deterministic gradient-corrected update rule at test time.

  We validate both routes on maximum clique, demonstrating that each recovers $O(n^2)$
  decoding while matching or improving on the $O(n^3)$ rerun state-of-the-art GNN different planted-clique regimes.
\end{abstract}

% ============================================================
%  INTRODUCTION
% ============================================================
\section{Introduction}

A wide class of algorithms for combinatorial optimization proceeds not by making a
single global decision, but by a sequence of \emph{residual decisions}: score the
current elements, prune or select one, update the problem state, and repeat. This
residual loop is the backbone of greedy peeling~\citep{feige2010}, spectral
rounding~\citep{alon1998}, local greedy algorithms, and branch-and-bound solvers. Its power lies in
adaptivity: each decision is informed by the current state of the problem rather than
a fixed global snapshot.

What makes this loop computationally tractable for classical algorithms is that the
quantities they maintain are \emph{additive}. Consider vertex degree as a concrete
example: when a vertex is removed from a graph, each remaining vertex's degree changes
by at most $1$, depending solely on whether it was adjacent to the removed vertex.
This local, fixed correction means the entire degree vector can be updated in $O(n)$
time, so a full residual decoder on a dense graph runs in $O(n^2)$ total. The same
pattern recurs across domains---marginal gains in submodular maximization, residual
correlations in spectral methods, variable beliefs in message passing---all amenable
to local additive updates after each decision~\citep{feige2010, dekel2014}. The
efficiency of the residual loop is, at its core, the efficiency of maintaining an
additive statistic.

\paragraph{GNNs as learned scorers.}
A natural ambition is to replace these hand-crafted statistics with a \emph{learned}
scoring rule that could capture richer structural patterns. Graph neural networks
(GNNs) have emerged as the leading candidate for this
role~\citep{karalias2020, min2022, garmendia2024}: trained end-to-end on problem
instances, they learn vertex rankings for problems such as \textsc{Max-Clique},
\textsc{Max-Cut}, \textsc{Max-Independent-Set}, and \textsc{TSP} without requiring
domain-specific feature engineering. Recent interpretability work has even shown that
GNNs trained on combinatorial tasks tend to rediscover classical statistics
spontaneously---the same degree-like ranking that classical peeling algorithms
use~\citep{shoham2026ltr, cappart2023}. This raises an optimistic picture: perhaps a
GNN can simply \emph{slot into} the residual loop in place of the hand-coded scorer.

\paragraph{The efficiency gap.}
The optimism, however, runs into a wall. GNNs are nonlinear, and nonlinearity breaks
the additive update property that makes the residual loop cheap. When a vertex is
removed, there is no $O(n)$ correction that reliably adjusts the remaining GNN scores;
the network must be \emph{rerun} on the current residual graph to produce fresh ones.
On a dense graph a single forward pass costs $O(n^2)$, and up to $O(n)$ reruns may be
needed, inflating total decoding time to $O(n^3)$---a full order of magnitude above
classical additive methods. The representational gains of neural scoring come at a
steep algorithmic price.

This raises the central question of this paper: \emph{can a GNN be designed to
  support efficient incremental decoding, recovering the $O(n^2)$ complexity of
  classical residual algorithms without giving up the expressive power of neural
scoring?}

\paragraph{Why ranking is not enough.}
A tempting first response is that a GNN which has learned to approximate degree should
automatically support additive updates, since degree itself is additive. We show this
intuition is wrong. The obstacle is not whether the model has learned a
degree-\emph{correlated} score, but whether it has learned a degree-\emph{affine}
one. Formally, a score $s_i = \varphi(d_i)$ supports exact additive deletion updates
if and only if $\varphi$ is affine (Theorem~\ref{thm:affine}). Any nonlinear
transformation of degree---no matter how tightly correlated with it---accumulates a
residual error at every pruning step. Worse, these errors compound: a single incorrect
removal sends the decoder onto a structurally different residual graph, and all
subsequent scores are computed on a different problem state, so the error at each
subsequent step is not a small perturbation but an arbitrary deviation
(Theorem~\ref{thm:propagation}). Our experiments confirm this sharply: a standard GNN
achieves near-perfect recovery under rerun-based decoding yet collapses to below
$0.35$ approximation under incremental decoding on the same instances. More tellingly,
adding \st{a training loss that imitates a rerun-based teacher---which we call \emph{Probable-Remover} (PR): at each step, the teacher reruns the GNN on the current residual graph and removes the lowest-scoring vertex; the student, running the GNN only once, is trained to predict the same removal sequence from incrementally updated scores---} \change{training to support the \emph{Probable-Remover} (PR) procedure---which at each step removes the lowest-scoring vertex and updates the remaining scores incrementally, applicable at both train and test time---} makes
the standard GNN \emph{worse} ($0.116$)---because without the right internal signal,
the imitation objective is unachievable and disrupts the static ranking that objective
training alone would have produced.

\paragraph{\change{Two routes to additive decoding.}}
The root cause is clear: the GNN lacks access to the additive marginal information
that classical residual algorithms maintain explicitly. We show two complementary
routes to restore such information, both grounded in the objective gradient.

\emph{Route 1 --- learned gradient updater (SGNN+GU/UPR).}
The first route keeps the GNN backbone standard, but gives the learned score updater
$U_\theta$ access to the objective gradient. After a vertex is removed, the updater
receives the current scores, the adjacency relation to the removed vertex, and the
current objective-gradient values. Since the gradient of a locally decomposable
objective has a local deletion correction, $U_\theta$ receives the information needed
to learn the residual score update. This route is flexible: the update rule is learned
rather than hard-coded.

\emph{Route 2 --- additive-channel GNN (ACGNN/ADD-LPR).}
The second route removes the learned updater entirely. ACGNN injects the objective
gradient
$
g_v=\partial \mathcal{L}/\partial p_v
$
inside the nonlinear message-passing backbone. After $L$ message-passing steps, the
additive coordinate is read only from the final hidden representation,
$
a_v = W_a h_v^{(L)}+b_a.
$
Thus the decoder is not given degree, soft degree, $(Ap)_v$, or a direct skip from
the gradient. The model must represent the additive coordinate through its nonlinear
embedding. We train this coordinate with marginal and residual-consistency losses so
that it behaves like an update-compatible state variable. At test time, decoder we denote ADD-LPR runs
the GNN once and updates $a_v$ deterministically using the local gradient correction
implied by the clique objective.

The two routes differ in what is learned. Route 1 learns the update rule through
$U_\theta$. Route 2 learns an additive coordinate inside the GNN representation and
uses an analytic residual update at test time. Both reduce dense-graph decoding from
$O(n^3)$ rerun decoding to $O(n^2)$ incremental decoding.

\change{One might ask whether the additive coordinate is simply a skip connection. It is
  not. A skip connection re-uses an opaque hidden state fixed at the end of the forward
  pass. ACGNN instead injects the objective gradient into message passing and learns
  $a_v$ only from the final nonlinear hidden state. ADD-LPR then updates this learned
  coordinate with the analytic clique-gradient correction; it is not a soft-degree skip
connection or a hand-coded $(Ap)_v$ feature.}

\paragraph{Contributions.}
On the theory side:

\medskip
\noindent\textbf{Theorem~\ref{thm:affine}.} Exact additive updates require affine
score encodings: the only functions compatible with a fixed additive correction after
each deletion are affine.

\noindent\textbf{Theorem~\ref{thm:nonlinear}.} Nonlinear GNN activations make affine
encodings unachievable in the raw score coordinate: bounded activations produce
bounded outputs, while any non-trivial affine function is unbounded. Ranking
supervision alone cannot rescue this, because it is blind to the additive scale that
exact residual updates require.

\noindent\textbf{Theorem~\ref{thm:propagation}.} Per-step errors compound: a single
wrong removal sends the decoder onto a structurally different residual graph, after
which subsequent errors are not small perturbations of the original trajectory but
may become arbitrary deviations.

\medskip
\noindent\change{\textbf{Theorem~\ref{thm:gu_positive}} (Gradient-updater
  representability). For locally decomposable objectives, the objective gradient has an
  exact local deletion rule: after removing vertex \(v\), each remaining marginal
  \(g_i\) changes by the contribution of the deleted local factor \((i,v)\). For
  quadratic objectives this correction is an affine local function,
  \[
    \Delta g_i(v)
    =
    -2\sum_r \alpha_r (M_r)_{iv}p_v .
  \]
  Therefore an updater \(U_\theta\) that receives gradient information can represent
  the residual correction needed by UPR. This is the positive result for SGNN+GU: the
gradient input exposes the additive signal required for local score updates.}

\medskip
\noindent\change{\textbf{Theorem~\ref{thm:acgnn_positive}} (Recoverability through
  nonlinear processing). Suppose the gradient values lie in a bounded interval and the
  activation used in the GNN backbone is continuous, strictly monotone, and
  non-saturated on that interval. Then the activation is injective there, so it does
  not erase the information in the gradient. A downstream readout from the final hidden
  state can recover affine functions of the gradient to arbitrary accuracy. Thus ACGNN
  can represent an additive coordinate \(a_i\approx -g_i\) even after the gradient has
  passed through nonlinear message-passing layers. This is a representability result,
  not a claim that the model automatically uses the signal without the marginal and
residual-consistency losses.}

\medskip
\noindent The first three results explain why standard GNNs fail at incremental
decoding; the last two give the positive routes: a learned gradient updater for
SGNN+GU and a recoverable additive coordinate for ACGNN.

\medskip
On the empirical side, we validate on planted clique with a principled
easy/medium/hard benchmark~\citep{shoham2026ltr}. The two routes give complementary
findings:

\begin{itemize}
  \item \change{\textbf{Route 1 (GU/UPR):} Feeding the gradient to $U_\theta$ is
      necessary and sufficient for learned incremental decoding. Without it, UPR
      collapses ($\sim$0.12 on Easy); with it, SGNN+GU achieves perfect
      Easy/Medium recovery. Notably, adding GU training to a standard GNN without the
      gradient signal makes performance \emph{worse}, confirming the gradient is
    necessary, not just helpful.}
  \item \change{\textbf{Route 2 (ACGNN/ADD-LPR):} The gradient-informed additive
      coordinate with consistency training achieves 1.000/1.000/0.886 on Easy/Medium/Hard
      under the fully deterministic ADD-LPR decoder---the strongest Hard result
      across all models---with no learned updater at test time. The consistency loss
    is key: objective training alone gives weaker Hard performance.}
  \item \change{Both routes reduce dense-graph decoding from $O(n^3)$ to $O(n^2)$.
      Route 1 is more flexible (any GNN + learned updater). Route 2 is more
    principled (deterministic, no test-time parameters beyond the GNN itself).}
\end{itemize}

\paragraph{Broader significance.}
While we use planted clique as our testbed, the challenge is general. Wherever a GNN
is inserted into a residual loop---peeling for independent set, variable elimination
in constraint satisfaction, decimation in belief propagation---the same tension
between representational power and incremental efficiency arises. \change{Our
  framework offers two complementary remedies: a learned gradient updater for
flexibility, and a gradient-informed additive coordinate for deterministic guarantees. Both} turn
any differentiable combinatorial objective into a source of cheap algorithmic state,
bridging the gap between neural scorers and the $O(n^2)$ efficiency of classical
additive algorithms.

\paragraph{Paper organization.}
Section~\ref{sec:related} reviews related work. Section~\ref{sec:method} develops the
methodology. Section~\ref{sec:theory} presents the theoretical analysis.
Section~\ref{sec:experiments} presents the empirical evaluation.
Section~\ref{sec:discussion} discusses limitations and future directions.

% ============================================================
%  RELATED WORK
% ============================================================
\section{Related Work}
\label{sec:related}

\change{\paragraph{Message-passing GNNs.}
  A \emph{message-passing neural network} (MPNN) maintains a hidden state
  $h_v^{(\ell)} \in \mathbb{R}^d$ at each vertex, updated over $L$ layers as
  \[
    h_v^{(\ell+1)} = \sigma\!\left(\mathbf{M}\!\left[h_v^{(\ell)},\;
    \mathrm{Agg}_{u \in \mathcal{N}(v)} h_u^{(\ell)}\right]\right),
  \]
  where $\mathbf{M}$ is a learnable MLP, $\mathrm{Agg}$ is a permutation-invariant
  aggregation, and $\sigma$ is a nonlinear activation~\citep{xu2019}. A per-vertex
  score $z_v = \mathbf{W}_{\mathrm{out}} h_v^{(L)}$ is read off after $L$ steps.
  Their expressive power is bounded by the 1-WL graph isomorphism
  test~\citep{xu2019, morris2019}, though our results concern an orthogonal
  limitation---\emph{update compatibility}---that persists even beyond the WL bound.
  MPNNs have found particular traction in combinatorial optimization, where their
permutation-invariant structure naturally fits graph-structured problems.}

\paragraph{\change{What GNNs learn on combinatorial problems.}}
\change{Understanding what MPNNs actually learn---rather than what they are capable of expressing---is a distinct and equally important question.} \citet{cappart2023} argue broadly that GNNs on NP-hard problems tend to
behave more like polynomial-time heuristics than like exact solvers.
\citet{shoham2026ltr} provide direct evidence for this on \textsc{Max-Clique}:
through PCA analysis of latent spaces, they show that GNNs consistently learn
degree-based vertex ranking across architectures, datasets, and even across domains
(the same model transfers to sparse PCA recovery). The concept-driven analysis
pipeline of~\citet{shoham2026sat} applies similar methods to SAT-solving GNNs,
identifying the algorithmic concepts that emerge from training. These findings
motivate our work directly: a GNN that has learned to approximate degree should, in
principle, support $O(n)$ residual updates, since degree itself is additive. We show
that in practice it does not, and that the reason is a precise algebraic
incompatibility between nonlinear score encodings and exact additive update rules\change{---which brings us to the question of how GNNs are used in practice for combinatorial optimization}.

\paragraph{GNNs for combinatorial optimization.}
An influential line of work frames
combinatorial problems as quadratic unconstrained binary optimization and trains GNNs
with differentiable relaxations of the
objective~\citep{schuetz2022, karalias2020, min2022}. \citet{karalias2020}
(Erd\H{o}s Goes Neural) introduced an unsupervised framework in which the GNN loss
is derived directly from the problem objective, and demonstrated it on
\textsc{Max-Clique} and graph partitioning. \citet{min2022} extended this to the
scattering-based architecture (SNN) we study here, obtaining strong results for
planted clique recovery. Beyond clique problems, similar unsupervised objectives have
been applied to \textsc{Max-Cut}, \textsc{Max-Independent-Set}, \textsc{TSP}, and
\textsc{Max-3-SAT}~\citep{schuetz2022, garmendia2024}. \change{All of these works
  insert the trained GNN into a decoding loop---the very setting where update
compatibility becomes critical.}

\paragraph{Architectural alignment with algorithms.}
\change{A complementary line} aligns the GNN architecture with the algebraic
structure of approximation algorithms. \citet{yau2023} (OptGNN) proved that
polynomial-sized message-passing networks can implement gradient descent on the
standard SDP relaxation for Max Constraint Satisfaction Problems, and that the
resulting architecture achieves UGC-optimal approximation guarantees in theory while
outperforming SDP solvers in practice. The key insight in OptGNN is that message
passing with learnable weights naturally implements projected gradient steps on a
low-rank SDP---the network's iterative aggregation is algebraically equivalent to
running an optimization algorithm. This is closely related to our own motivation: we
seek to align message passing with the linear marginal information that cheap
residual updates require. The difference is that OptGNN aligns GNNs with a convex
relaxation to obtain approximation guarantees, while our model-aware GNN aligns with
the local gradient of the \emph{learned} objective to enable efficient residual
decoding\change{---a distinction that connects directly to the role of gradient signals in neural solvers more broadly}.

\paragraph{Gradient and marginal signals in neural solvers.}
A recurring theme in neural combinatorial optimization is that objective-aware
signals improve solutions. \citet{schuetz2022} uses physics-inspired loss functions
derived from the problem Hamiltonian, effectively baking first-order information into
training. The recurrent feature update approach of QRF-GNN~\citep{gaile2024} feeds
intermediate GNN predictions back as dynamic node features, allowing the network to
move closer to the objective minimum across iterations; this is structurally similar
to our gradient channel, but targets convergence within a single inference pass
rather than residual update compatibility across sequential pruning steps. The
unrolled GNN of~\citet{hadou2024} explicitly unrolls a dual ascent algorithm into
coupled primal and dual networks, so that each layer computes a meaningful
algorithmic step. Algorithm unrolling more broadly~\citep{monga2021} provides a
systematic design principle: embed known iterative solvers into neural layers and
make their parameters learnable. Our self-aware architecture differs in that we do
not unroll a fixed algorithm---we expose the gradient of the differentiable objective
as an internal signal and let the model learn how to use it, which is more flexible
when the correct update rule is unknown\change{. This gives rise to two complementary
  approaches explored in this paper: a \emph{learned} gradient updater ($U_\theta$-GU),
  closer in spirit to algorithm unrolling in that it trains a network to approximate an
  iterative process; and a \emph{deterministic} gradient-corrected additive decoder
  (ACGNN/ADD-LPR),
  closer to classical algorithm design in that the update rule is analytically derived
from the objective and requires no learned components at test time.}

\paragraph{Imitation and self-supervised decoding.}
Our \change{PR} training regime is related to imitation learning approaches in
combinatorial optimization~\citep{liu2021}, in which a GNN is trained to replicate
the decisions of a stronger algorithm. The novelty in our setting is that the teacher
and student share weights: the student observes the teacher's rerun-based removal
sequence and is trained to predict it using only incrementally updated scores, a form
of self-imitation. This is also related to learning-to-search
methods~\citep{bengio2021} in which a policy is trained to follow expert
trajectories, but adapted to the specific problem of approximating adaptive residual
decoding with a single network pass.

% ============================================================
%  METHODOLOGY
% ============================================================
\section{Methodology}
\label{sec:method}

We evaluate four model families, chosen to isolate where additive residual
information enters the system:
\[
  \text{SGNN},\qquad
  \text{SGNN+U},\qquad
  \text{SGNN+GU},\qquad
  \text{ACGNN}.
\]
SGNN is a standard message-passing GNN. SGNN+U adds a learned updater $U_\theta$ but
does not give it gradient information. SGNN+GU gives the updater access to the
objective gradient. ACGNN instead injects the objective gradient into the nonlinear
GNN backbone and learns a final additive coordinate used by a deterministic decoder.

The corresponding decoders are RERUN, LPR, UPR, and ADD-LPR. RERUN recomputes GNN
scores after every deletion and costs $O(n^3)$ on dense graphs. LPR runs the GNN once
and prunes using fixed scores. UPR runs the GNN once and then updates scores using
$U_\theta$. ADD-LPR runs ACGNN once and then applies an analytic gradient-correction
update to the learned additive coordinate.

\subsection{Problem Setting}
\label{sec:problem}

We evaluate on the planted clique problem on $G(n,\tfrac{1}{2},k)$ random graphs.
Each instance is generated by first sampling an Erd\H{o}s--R\'enyi graph
$G(n,\tfrac{1}{2})$, then planting a clique $C^\star$ of size $k$ by connecting $k$
uniformly chosen vertices. The goal is to recover a large valid clique, ideally
$C^\star$.

\paragraph{Difficulty regimes.}
Following~\citet{shoham2026ltr}, we partition instances into three regimes based on
known theoretical thresholds for planted clique recovery:
\begin{itemize}
  \item \textbf{Easy} ($k \in [62, 100]$ for $n=1000$): degree-based top-$k$
    heuristics recover the planted clique with high probability~\citep{kucera1995}.
  \item \textbf{Medium} ($k \in [36, 61]$): recovery requires iterative residual
    methods such as degree-based LPR or spectral
    algorithms~\citep{alon1998, feige2010}.
  \item \textbf{Hard} ($k \in [20, 35]$): under the planted clique hardness
    conjecture, no polynomial-time algorithm is known to
    succeed~\citep{manurangsi2021}.
\end{itemize}

\paragraph{Evaluation metrics.}
We report three metrics on the recovered clique $C_{\text{found}}$:
\begin{align*}
  \text{Approximation} &= |C_{\text{found}}| / |C^\star|, \\
  \text{Overlap}       &= |C_{\text{found}} \cap C^\star| / |C^\star|, \\
  \text{Exact}         &= \mathbf{1}[C_{\text{found}} = C^\star].
\end{align*}
All experiments use $n = 1000$, averaged over 1000 test instances per regime.

\subsection{Model Architectures}
\label{sec:architectures}

All models are message-passing GNNs with $L=10$ shared message-passing steps, hidden
dimension $64$, and $\sigma=\tanh$ in the reported implementation. The input feature
is the same for all vertices and contains no degree, clique-membership, or other
hand-designed structural feature.

\paragraph{Standard GNN (SGNN).}
SGNN is a standard message-passing GNN. At layer $\ell$, each vertex aggregates
neighbor representations
\[
  q_v^{(\ell)}
  =
  \mathrm{Agg}_{u\in\mathcal{N}(v)} h_u^{(\ell)}
\]
and updates its hidden state by
\[
  h_v^{(\ell+1)}
  =
  \sigma\!\left(
    M_{\mathrm{upd}}\!\left[
      h_v^{(\ell)},q_v^{(\ell)}
    \right]
  \right).
\]
The final score is
\[
  s_v = W_{\mathrm{out}} h_v^{(L)} + b_{\mathrm{out}} .
\]
This score is used by one-pass, LPR, and rerun decoding.

\paragraph{SGNN+U and SGNN+GU.}
SGNN+U and SGNN+GU use the same SGNN backbone, but add a learned score updater
$U_\theta$ for UPR decoding. After vertex $v^*$ is removed, each remaining score is
updated as
\[
  s_i
  \leftarrow
  s_i
  +
  U_\theta(s_i,s_{v^*},A_{iv^*},\gamma_i,\gamma_{v^*}) .
\]
For SGNN+U, the gradient channels are zero:
\[
  \gamma_i=\gamma_{v^*}=0.
\]
For SGNN+GU, the updater receives the current objective-gradient values:
\[
  \gamma_i =
  g_i
  =
  \frac{\partial \mathcal{L}_{\mathrm{clique}}}{\partial p_i}.
\]
Thus SGNN+GU tests whether gradient information is useful specifically at the
residual update stage, without changing the GNN backbone.

\paragraph{Additive-Channel GNN (ACGNN).}
ACGNN removes the learned updater and instead trains the GNN backbone to expose an
additive coordinate. At internal layer $\ell$, the current hidden states are scored as
\[
  z_v^{(\ell)}
  =
  W_{\mathrm{out}} h_v^{(\ell)} + b_{\mathrm{out}},
  \qquad
  p_v^{(\ell)}
  =
  \operatorname{sigmoid}(z_v^{(\ell)}).
\]
The model then computes the objective-gradient marginal
\[
  g_v^{(\ell)}
  =
  \frac{\partial
  \mathcal{L}_{\mathrm{clique}}(p^{(\ell)};A,k)}
  {\partial p_v^{(\ell)}} .
\]
This gradient is injected into the nonlinear message-passing update:
\[
  h_v^{(\ell+1)}
  =
  \sigma\!\left(
    M_{\mathrm{upd}}\!\left[
      h_v^{(\ell)},q_v^{(\ell)},g_v^{(\ell)}
    \right]
  \right).
\]
After $L$ steps, ACGNN outputs both the standard score
\[
  s_v = W_{\mathrm{out}}h_v^{(L)}+b_{\mathrm{out}},
\]
and an additive coordinate
\[
  a_v = W_a h_v^{(L)}+b_a .
\]
There is no direct path from $g_v$, $(Ap)_v$, degree, or any hand-designed statistic
to $a_v$. The additive coordinate must be represented in the final nonlinear hidden
state.

\subsection{Decoding Procedures}
\label{sec:decoding}

\change{We evaluate three decoders. All start from an initial GNN forward pass and
  iteratively remove the lowest-scoring vertex until the remaining set forms a clique,
followed by an add-back pass.}

\paragraph{Rerun decoder (RERUN).}
After each removal, the full GNN is rerun on the current residual graph to produce
fresh scores. Adaptive baseline; costs $O(n^3)$ on dense graphs.

\paragraph{PR decoder (PR / UPR).}
The GNN runs once. After each removal of vertex $v^*$, a learned node-wise updater
$U_\theta$ adjusts remaining scores:
\begin{equation}
  s_i \leftarrow s_i + U_\theta(s_i,\, s_{v^*},\, A_{i v^*},\, g_i,\, g_{v^*}),
  \qquad \forall i \in V(G_R),
  \label{eq:updater}
\end{equation}
\change{where $g_i, g_{v^*}$ are the objective-gradient values (zero for SGNN).
  $U_\theta$ is a separate lightweight MLP, independent of the GNN backbone, trained
jointly with the rest of the model.} Since $U_\theta$ is applied independently to
each remaining vertex, each removal costs $O(n)$, giving $O(n^2)$ total.

\change{\paragraph{Why $g_i$ in $U_\theta$ enables additive updates.}
  The correct removal ordering is $\arg\min_i g_i$. By
  Theorem~\ref{thm:gu_positive}, $\Delta g_i$ under each deletion is a linear function
  of inputs already available to $U_\theta$ when $g_i$ is provided. A linear first
  layer can therefore represent $\Delta g_i$ exactly, before any nonlinear activation
  of $U_\theta$ fires. Without $g_i$, recovering $\Delta g_i$ from $s_{v^*}=\sigma(z_{v^*})$
  requires inverting $\sigma$, which by Theorem~\ref{thm:nonlinear} cannot be achieved
with a fixed additive rule. A formal proof appears in Appendix~\ref{app:rectification}.}

\paragraph{Additive decoder (ADD-LPR).}
ADD-LPR is used only with ACGNN. The GNN runs once on the original graph and returns
the additive coordinate $a_i$ and probabilities $p_i$. The decoder then repeatedly
removes the lowest-scoring vertex under $a_i$. After removing vertex $v^*$, it updates
the remaining additive scores using the analytic correction to the clique gradient:
\begin{equation}
  a_i
  \leftarrow
  a_i-\Delta g_i(v^*),
  \qquad i\in V(G_R),
  \label{eq:additive_update}
\end{equation}
where
\[
  \Delta g_i(v^*)
  =
  g_i(G_R\setminus\{v^*\})-g_i(G_R).
\]
For the clique objective, this correction is
\[
  \Delta g_i(v^*)
  =
  \frac{2A_{iv^*}p_{v^*}}{k(k-1)}
  -
  \frac{2B_{iv^*}p_{v^*}}{k(n-k)}
  -
  \frac{20p_{v^*}}{n^2}.
\]
The sign follows our convention that ACGNN is trained to satisfy $a_i\approx -g_i$.
Thus, when $g_i$ changes by $\Delta g_i$, the negative-gradient coordinate changes by
$-\Delta g_i$.

\begin{algorithm}[t]
  \caption{Additive Decoder (ADD-LPR)}
  \label{alg:add_lpr}
  \begin{algorithmic}[1]
    \State Run ACGNN once on $G$ to obtain $a_i$ and $p_i$ for all vertices.
    \State $C\leftarrow V(G)$;\quad $H\leftarrow []$.
    \While{$C$ does not form a clique}
    \State $v^*\leftarrow \arg\min_{i\in C} a_i$.
    \State Append $v^*$ to $H$.
    \State $C\leftarrow C\setminus\{v^*\}$.
    \State Compute $\Delta g_i(v^*)$ for all $i\in C$.
    \State $a_i\leftarrow a_i-\Delta g_i(v^*)$ for all $i\in C$.
    \EndWhile
    \For{$v$ from last to first in $H$}
    \If{$C\cup\{v\}$ is a clique}
    \State $C\leftarrow C\cup\{v\}$.
    \EndIf
    \EndFor
    \State \Return $C$.
  \end{algorithmic}
\end{algorithm}

\subsection{Training Regimes}
\label{sec:training}

Both regimes train on Medium-difficulty instances. They share the clique objective
but differ in whether an additional imitation signal is used to explicitly supervise
the \change{PR} decoder.

\paragraph{Objective training (OBJ).}
The model is trained with the unsupervised clique objective alone:
\begin{equation}
  \mathcal{L}_{\mathrm{train}} = \mathcal{L}_{\mathrm{clique}}.
\end{equation}
$U_\theta$ is learned only through this objective, with no direct signal about the
removal sequence. This tests whether the architecture alone is sufficient to induce
update-compatible representations.

\paragraph{Probable-Remover training (PR).}
PR training adds an imitation loss that directly supervises the \change{PR} decoder.
A rerun-based \emph{teacher} at each step $t$ reruns the GNN on the current residual
graph $G_t$ and removes the lowest-scoring vertex $v^t_{\mathrm{teacher}}$. A
\emph{student} runs the GNN only once, then tracks the teacher's removal sequence
using incrementally updated scores $\tilde{s}^{(t)}$, with $U_\theta$ applied after
each step. The imitation loss penalises disagreement between student and teacher:
\begin{equation}
  \mathcal{L}_{\mathrm{PR}} = \frac{1}{T} \sum_{t=0}^{T-1}
  \mathrm{CE}_{\mathrm{softmax}}\!\left(-\tilde{s}^{(t)}_{V(G_t)},\,
  v^t_{\mathrm{teacher}}\right),
\end{equation}
where the minus sign reflects that lower scores correspond to removed vertices. The
total loss is:
\begin{equation}
  \mathcal{L}_{\mathrm{train}} =
  \mathcal{L}_{\mathrm{clique}} + \lambda_{\mathrm{PR}}\, \mathcal{L}_{\mathrm{PR}},
\end{equation}
with $\lambda_{\mathrm{PR}} = 1$. Algorithm~\ref{alg:training} summarizes one PR
training step.

\begin{algorithm}[t]
  \caption{PR Training Step}
  \label{alg:training}
  \begin{algorithmic}[1]
    \State Sample a Medium-difficulty instance $(G, k)$
    \State $\tilde{s} \leftarrow \mathrm{GNN}(G)$ \Comment{Student: single forward pass}
    \State $G_R \leftarrow G$;\quad $\mathcal{L}_{\mathrm{PR}} \leftarrow 0$
    \For{$t = 1, \ldots, T$}
    \State $s^T \leftarrow \mathrm{GNN}(G_R)$ \Comment{Teacher: rerun on residual graph}
    \State $v \leftarrow \arg\min_{u \in V(G_R)} s^T_u$ \Comment{Teacher removes lowest}
    \State $\mathcal{L}_{\mathrm{PR}} \mathrel{+}=
    \mathrm{CE}_{\mathrm{softmax}}(-\tilde{s}_{V(G_R)},\, v)$
    \Comment{Student must predict same removal}
    \State $G_R \leftarrow G_R \setminus \{v\}$
    \State $\tilde{s}_i \leftarrow \tilde{s}_i +
    U_\theta(\tilde{s}_i, \tilde{s}_v, A_{iv}, g_i, g_v)$
    \quad for all $i \in V(G_R)$ \Comment{Student updates scores incrementally}
    \EndFor
    \State $\mathcal{L}_{\mathrm{train}} =
    \mathcal{L}_{\mathrm{clique}} + \mathcal{L}_{\mathrm{PR}} / T$
    \State Optimize $\mathcal{L}_{\mathrm{train}}$
  \end{algorithmic}
\end{algorithm}

\paragraph{Additive-coordinate training (AC).}
AC training is used only for ACGNN. Its purpose is to make the final coordinate
$a_i=W_a h_i^{(L)}+b_a$ behave as an update-compatible residual state. Since lower
scores are removed, we train $a_i$ to approximate the negative objective gradient:
\[
  a_i(G) \approx -g_i(G).
\]
The marginal alignment loss is
\[
  \mathcal{L}_{\mathrm{marginal}}
  =
  \frac{1}{|V|}
  \sum_{i\in V}
  \left(a_i(G)+g_i(G)\right)^2 .
\]
We also enforce residual consistency. During training only, we sample vertices $v$,
rerun ACGNN on the residual graph $G\setminus\{v\}$, and require the new additive
coordinate to agree with the analytic gradient correction:
\[
  \mathcal{L}_{\mathrm{res}}
  =
  \mathbb{E}_{v\sim V}
  \frac{1}{|V|-1}
  \sum_{i\neq v}
  \left(
    a_i(G\setminus\{v\})
    -
    \left[
      a_i(G)-\Delta g_i(v)
    \right]
  \right)^2 .
\]
The total training objective is
\[
  \mathcal{L}_{\mathrm{train}}
  =
  \mathcal{L}_{\mathrm{clique}}
  +
  \lambda_m\mathcal{L}_{\mathrm{marginal}}
  +
  \lambda_r\mathcal{L}_{\mathrm{res}} ,
\]
with $\lambda_m=\lambda_r=1$ in the reported experiments. This loss does not provide
degree, soft degree, or planted labels. It supervises the final hidden representation
to expose the objective-gradient marginal as an additive coordinate.

The evaluated model/training/decoder combinations are:
\[
  \begin{array}{lll}
    \textbf{SGNN}     & \text{OBJ}      & \text{RERUN, LPR},\\
    \textbf{SGNN+U}   & \text{OBJ}      & \text{RERUN, LPR, UPR},\\
    \textbf{SGNN+GU}  & \text{OBJ, PR}  & \text{RERUN, LPR, UPR},\\
    \textbf{ACGNN}    & \text{OBJ, AC}  & \text{RERUN, LPR, ADD-LPR}.
  \end{array}
\]

% ============================================================
%  THEORETICAL ANALYSIS
% ============================================================
\section{Theoretical Analysis}
\label{sec:theory}

We establish four results. The first three characterize why standard GNNs cannot
support cheap incremental updates; the last two give the positive routes: a learned gradient updater for SGNN+GU and a recoverable additive coordinate for ACGNN. Throughout, $d_i$ denotes an integer-valued additive statistic such as residual
degree.

\begin{theorem}[Affine encoding is necessary for exact additive updates]
  \label{thm:affine}
  Let $d_i \in \mathbb{Z}_{\geq 0}$ be an additive statistic that decreases by
  exactly $1$ when a neighbor is deleted. Suppose the score $s_i = \varphi(d_i)$ is
  updated by a fixed additive rule $s_i \leftarrow s_i - c$ for some constant $c$
  whenever a neighbor is deleted. If this update is exact for all values of $d_i$ in
  the relevant range, then $\varphi$ must be affine: $\varphi(d) = cd + b$ for some
  $b \in \mathbb{R}$.
\end{theorem}
\begin{proof}
  Exactness requires $\varphi(d - 1) = \varphi(d) - c$ for all $d$, so the first
  discrete difference of $\varphi$ is the constant $c$. Any such function satisfies
  $\varphi(d) = cd + b$ for $b = \varphi(0)$.
\end{proof}

\begin{theorem}[Nonlinear GNNs cannot implement exact additive updates]
  \label{thm:nonlinear}
  Let $f_\theta : \mathcal{G} \to \mathbb{R}^n$ be a message-passing GNN whose
  update MLPs contain at least one bounded non-polynomial activation $\sigma$
  (e.g.\ $\tanh$, sigmoid). Then for any fixed constant $c$, there exists a graph
  $G$ and adjacent vertices $v, u$ such that $f_\theta(G)_v - c \neq
  f_\theta(G \setminus \{u\})_v$, regardless of the parameter choice $\theta$.
\end{theorem}
\begin{proof}
  By Theorem~\ref{thm:affine}, exact updates require $f_\theta(G)_v$ to be affine
  in $d_v$ for every $G$. Consider star graphs $K_{1,d}$ with $v$ as hub. With
  constant initial features, $f_\theta(K_{1,d})_v =: \varphi_\theta(d)$ depends
  only on $d$, and each message-passing step applies $\sigma$ to a linear function
  of $d$. For $\sigma = \tanh$, $\varphi_\theta$ is bounded, while any affine
  function with $c \neq 0$ is unbounded; $c = 0$ gives a trivial update that changes
  no scores. In either case some $d$ satisfies $\varphi_\theta(d-1) \neq
  \varphi_\theta(d) - c$, incurring nonzero error.
\end{proof}

\begin{remark}
  Theorems~\ref{thm:affine} and~\ref{thm:nonlinear} establish an architectural
  impossibility: no nonlinear GNN can support exact additive updates. A complementary
  training-side observation reinforces this: even if one tried to enforce affine
  structure through supervision, ranking loss cannot do so. Any strictly increasing
  reparametrisation $\varphi(d_i)$ of degree satisfies the same pairwise ranking
  constraints as $d_i$ itself, so ranking supervision is blind to the additive scale
  that exact updates require. The impossibility is therefore both architectural and
  statistical.
\end{remark}

\begin{theorem}[Local errors cascade into trajectory divergence]
  \label{thm:propagation}
  Consider iterative pruning that removes the vertex with the lowest true score $s_i$
  at each step. If instead estimated scores $\hat{s}_i = s_i + \varepsilon_i$ are
  used, and at some step $t$ there exist $i \neq j$ with
  \begin{equation}
    s_i < s_j \quad \text{but} \quad \hat{s}_i > \hat{s}_j,
    \label{eq:flip}
  \end{equation}
  then the estimated procedure removes $j$ while the exact procedure removes $i$.
  The resulting residual graphs differ, and all subsequent errors are not bounded by
  $\max_i |\varepsilon_i|$: they equal the difference between scores computed on
  structurally different graphs, which can be arbitrarily large.
\end{theorem}
\begin{proof}
  Condition~\eqref{eq:flip} requires $|\varepsilon_i - \varepsilon_j| > |s_i -
  s_j|$, so the estimated procedure removes $j$ and the exact procedure removes $i$,
  giving $G^t = G^{t-1} \setminus \{i\}$ and $\hat{G}^t = G^{t-1} \setminus \{j\}$.
  At step $t+1$, scores are computed on structurally different graphs: vertex $i$
  is present in $\hat{G}^t$ but absent from $G^t$, so neighborhoods differ. The
  deviation at step $t+1$ is not $\varepsilon_v$ but the difference between GNN
  outputs on different graphs, which can be arbitrarily large. The error compounds
  rather than averages out.
\end{proof}

\begin{theorem}[Gradient-updater representability]
  \label{thm:gu_positive}
  Let
  \[
    \mathcal{L}(p;G)
    =
    \sum_i \phi_i(p_i)
    +
    \sum_{(i,j)\in F(G)}
    \psi_{ij}(p_i,p_j,A_{ij})
  \]
  be a locally decomposable objective, and let
  \[
    g_i(G)=\frac{\partial \mathcal{L}(p;G)}{\partial p_i}.
  \]
  After deleting vertex \(v\), for every remaining vertex \(i\neq v\),
  \[
    g_i(G\setminus \{v\})
    =
    g_i(G)
    -
    \frac{\partial \psi_{iv}(p_i,p_v,A_{iv})}{\partial p_i},
  \]
  with the convention that the correction is zero when \((i,v)\notin F(G)\).
  Therefore the residual correction
  \[
    \Delta g_i(v)
    =
    g_i(G\setminus \{v\})-g_i(G)
  \]
  is a local function of the deleted factor involving \(i\) and \(v\), rather than a
  function requiring a global GNN rerun.

  For quadratic objectives of the form
  \[
    \mathcal{L}(p;G)
    =
    \sum_r \alpha_r p^\top M_r p
    +
    \text{terms independent of the deleted edge factors},
  \]
  the correction is
  \[
    \Delta g_i(v)
    =
    -2\sum_r \alpha_r (M_r)_{iv}p_v .
  \]
  Hence an updater \(U_\theta\) that receives gradient information can represent the
  local residual correction required by UPR. This establishes the positive result for
  SGNN+GU: the gradient channel exposes the additive signal needed to update residual
  scores locally.
\end{theorem}

\begin{proof}
  By local decomposability,
  \[
    g_i(G)
    =
    \phi_i'(p_i)
    +
    \sum_{j:(i,j)\in F(G)}
    \frac{\partial \psi_{ij}(p_i,p_j,A_{ij})}{\partial p_i}.
  \]
  Deleting \(v\) removes exactly the factor \(\psi_{iv}\) from this sum and leaves all
  other local factors involving \(i\) unchanged. Therefore,
  \[
    g_i(G\setminus \{v\})
    =
    g_i(G)
    -
    \frac{\partial \psi_{iv}(p_i,p_v,A_{iv})}{\partial p_i}.
  \]
  This proves the local deletion rule.

  For a quadratic term \(p^\top M_r p\), using the symmetric contribution, the
  \(i\)-th gradient coordinate is
  \[
    \nabla_i(p^\top M_r p)=2(M_rp)_i.
  \]
  Removing \(v\) removes the \(v\)-coordinate contribution from this marginal, giving
  \[
    \Delta g_i(v)
    =
    -2\sum_r \alpha_r (M_r)_{iv}p_v .
  \]
  Thus the required correction is determined by local quantities around the deleted
  vertex and can be represented by a learned updater that receives gradient features.
\end{proof}

\begin{theorem}[Recoverability through nonlinear processing]
  \label{thm:acgnn_positive}
  Let the gradient values \(g_i\) lie in a bounded interval \(I\subset\mathbb{R}\).
  Assume the activation \(\sigma\) used in the ACGNN backbone is continuous, strictly
  monotone, and non-saturated on \(I\). Then \(\sigma\) is injective on \(I\), and
  there exists a continuous inverse \(\sigma^{-1}\) on \(\sigma(I)\). Consequently,
  for every affine function \(\ell(g)=\alpha g+\beta\) and every \(\varepsilon>0\),
  there exists an MLP readout \(R_\theta\) such that
  \[
    \sup_{g\in I}
    \left|
    R_\theta(\sigma(g))-\ell(g)
    \right|
    <\varepsilon .
  \]
  In particular, ACGNN can represent an additive coordinate \(a_i\approx -g_i\) after
  the gradient has passed through nonlinear message-passing layers.
\end{theorem}

\begin{proof}
  Since \(\sigma\) is continuous and strictly monotone on \(I\), it is one-to-one on
  \(I\). Therefore it has a well-defined inverse \(\sigma^{-1}\) on the image
  \(\sigma(I)\). Because \(I\) is compact and \(\sigma\) is continuous and strictly
  monotone, \(\sigma(I)\) is compact and \(\sigma^{-1}\) is continuous.

  For any affine target \(\ell(g)=\alpha g+\beta\), define
  \[
    f(z)=\ell(\sigma^{-1}(z)),
    \qquad z\in \sigma(I).
  \]
  The function \(f\) is continuous on the compact set \(\sigma(I)\). By the universal
  approximation theorem, an MLP readout \(R_\theta\) can approximate \(f\) uniformly
  on \(\sigma(I)\). Hence, for every \(\varepsilon>0\), there exists \(R_\theta\) such
  that
  \[
    \sup_{z\in \sigma(I)}
    |R_\theta(z)-f(z)|<\varepsilon .
  \]
  Substituting \(z=\sigma(g)\) gives
  \[
    \sup_{g\in I}
    \left|
    R_\theta(\sigma(g))-\ell(g)
    \right|
    <\varepsilon .
  \]
  Taking \(\ell(g)=-g\) proves that the readout can recover the additive coordinate
  \(a_i\approx -g_i\). Thus monotone non-saturated nonlinear processing does not erase
  the gradient information needed for additive residual decoding.
\end{proof}

\begin{corollary}[Clique-gradient correction]
  \label{cor:clique_gradient}
  For the clique objective
  \[
    \mathcal{L}_{\mathrm{clique}}(p;A,k)
    =
    -\frac{p^\top A p}{k(k-1)}
    +
    \frac{p^\top B p}{k(n-k)}
    +
    10\left(\frac{\sum_i p_i-k}{n}\right)^2 ,
  \]
  where \(B\) is the non-edge matrix with zero diagonal, the gradient coordinate is
  \[
    g_i
    =
    -\frac{2(Ap)_i}{k(k-1)}
    +
    \frac{2(Bp)_i}{k(n-k)}
    +
    \frac{20(\sum_j p_j-k)}{n^2}.
  \]
  After deleting vertex \(v\), the local correction is
  \[
    \Delta g_i(v)
    =
    g_i(G\setminus\{v\})-g_i(G)
    =
    \frac{2A_{iv}p_v}{k(k-1)}
    -
    \frac{2B_{iv}p_v}{k(n-k)}
    -
    \frac{20p_v}{n^2}.
  \]
  Therefore, if ACGNN is trained so that \(a_i\approx -g_i\), the deterministic
  ADD-LPR update is
  \[
    a_i \leftarrow a_i-\Delta g_i(v).
  \]
\end{corollary}

\begin{remark}
  The gradient does not force the model to be additive. Rather, it supplies additive
  residual information. The GU updater can use this information directly to represent
  local corrections, while ACGNN can preserve it through monotone non-saturated
  nonlinear layers and expose it through a supervised additive readout. The marginal
  and residual-consistency losses in Section~\ref{sec:training} are what encourage
  the model to actually expose this representable signal as \(a_i\approx -g_i\).
\end{remark}

% ============================================================
%  EXPERIMENTS
% ============================================================
\section{Experiments}
\label{sec:experiments}

\subsection{Experimental Setup}
All models use $L=10$ shared message-passing steps, hidden dimension $64$, and
$\sigma = \tanh$. Models are trained for 100 epochs with batch size 32 on
Medium-difficulty planted clique instances (the choice of training difficulty has
minimal effect on the learned concepts, following~\citet{shoham2026ltr}).
\change{All experiments use $n=1000$ and results are averaged over 1000 test
instances per regime.} All experiments were run on a desktop machine equipped with an AMD Ryzen 9
5950X 16-core processor (32 logical cores, 3.4\,GHz), 64\,GB RAM, and an NVIDIA RTX
3080 GPU, running Windows 11 Pro. Code was implemented in Python 3.11 with PyTorch 2.2.

\subsection{Ablation: Gradient in $U_\theta$ vs Gradient in GNN Backbone}
\label{sec:ablation}

\change{Before presenting the full results, we isolate the key design decision: does
  the gradient need to go into $U_\theta$, into the GNN backbone, or both? We run four
  variants under the PR decoder, averaging across Easy/Medium/Hard regimes:

  \begin{table}[t]
    \centering
    \caption{\change{Ablation of gradient placement under the PR decoder. G = gradient
        in GNN backbone; U = gradient in $U_\theta$. Results are mean approximation ratio
        averaged over Easy/Medium/Hard regimes across three training regimes
    (OBJ/imitation/PR).}}
    \label{tab:ablation}
    \begin{tabular}{lcccccc}
      \toprule
      Model & G in GNN & G in $U_\theta$
      & \multicolumn{3}{c}{Approx (Easy / Medium / Hard)} \\
      \midrule
      SGNN\_U  & \xmark & \xmark & 0.467 & 0.293 & 0.267 \\
      SGNN\_GU & \xmark & \cmark & \textbf{0.954} & \textbf{0.947} & \textbf{0.867} \\
      MAGNN\_U & \cmark & \xmark & 0.243 & 0.003 & 0.000 \\
      MAGNN\_GU& \cmark & \cmark & 0.900 & 0.841 & 0.833 \\
      \bottomrule
    \end{tabular}
  \end{table}

  Three conclusions follow directly from Table~\ref{tab:ablation}. First, gradient in
  $U_\theta$ is what matters: SGNN\_GU (gradient only in $U_\theta$, not in the
  backbone) substantially outperforms all other variants including MAGNN\_U (gradient
  only in the backbone). Second, gradient in the GNN backbone alone is actively harmful:
  MAGNN\_U performs \emph{worse} than SGNN\_U, the baseline with no gradient anywhere.
  The backbone uses the gradient to produce better scores on the original graph, but
  those scores go stale immediately after the first deletion and provide no update
  signal. Third, combining both (MAGNN\_GU) is strong but does not dominate SGNN\_GU,
  confirming that the backbone gradient is secondary.

  This directly validates the theoretical argument in Section~\ref{sec:decoding}: the
  correct score update $\Delta g_i$ is a linear function of $U_\theta$'s inputs when
  $g_i$ is available, so $U_\theta$ can represent it exactly. Without $g_i$ in
$U_\theta$, no such representation exists regardless of what the backbone does.}
\label{sec:main_results}

Table~\ref{tab:main} reports planted clique recovery for all four model variants:
SGNN and MAGNN, each trained with either objective-only or self-predictor training,
and evaluated under both the \change{RERUN} and \change{PR} decoders. This
$2 \times 2 \times 2$ factorial design lets us isolate the contribution of each
component cleanly.

\begin{table}[t]
  \centering
  \caption{Planted clique recovery on $n=1000$.
    Approx $= |C_{\text{found}}|/|C^\star|$;
    Overlap $= |C_{\text{found}} \cap C^\star|/|C^\star|$;
    Exact $=$ fraction of instances with perfect recovery.
  Best incremental result in each regime is \textbf{bold}.}
  \label{tab:main}
  \resizebox{\textwidth}{!}{%
    \begin{tabular}{llccccccccc}
      \toprule
      & & \multicolumn{3}{c}{Easy} & \multicolumn{3}{c}{Medium}
      & \multicolumn{3}{c}{Hard} \\
      \cmidrule(lr){3-5}\cmidrule(lr){6-8}\cmidrule(lr){9-11}
      Model & Decoder & Approx & Overlap & Exact
      & Approx & Overlap & Exact
      & Approx & Overlap & Exact \\
      \midrule
      \multirow{2}{*}{SGNN + Obj}
      & Rerun     & 1.000 & 1.000 & 1.000
      & 1.000 & 1.000 & 1.000
      & 0.621 & 0.342 & --- \\
      & Increment & 0.354 & 0.282 & 0.200
      & 0.285 & 0.132 & 0.000
      & 0.390 & 0.023 & --- \\
      \midrule
      \multirow{2}{*}{MAGNN + Obj}
      & Rerun     & 1.000 & 1.000 & 1.000
      & 1.000 & 1.000 & 1.000
      & 0.607 & 0.339 & --- \\
      & Increment & 0.928 & 0.925 & 0.800
      & 0.772 & 0.740 & 0.600
      & 0.370 & 0.114 & --- \\
      \midrule
      \multirow{2}{*}{SGNN + Self}
      & Rerun     & 1.000 & 1.000 & 1.000
      & 1.000 & 1.000 & 1.000
      & 0.806 & 0.698 & --- \\
      & Increment & 0.116 & 0.003 & 0.000
      & 0.195 & 0.004 & 0.000
      & 0.373 & 0.008 & --- \\
      \midrule
      \multirow{2}{*}{MAGNN + Self}
      & Rerun     & 1.000 & 1.000 & 1.000
      & 0.698 & 0.600 & 0.600
      & 0.390 & 0.031 & --- \\
      & Increment & \textbf{1.000} & \textbf{1.000} & \textbf{1.000}
      & \textbf{1.000} & \textbf{1.000} & \textbf{1.000}
      & \textbf{0.689} & \textbf{0.427} & --- \\
      \bottomrule
    \end{tabular}%
  }
\end{table}

\paragraph{Rerun decoding is strong but expensive.}
Under rerun decoding, all models achieve perfect recovery on Easy and Medium. SGNN +
Self even reaches $0.806$ on Hard. This confirms the learning-to-rank phenomenon
of~\citet{shoham2026ltr}: with unrestricted access to fresh GNN scores after every
removal, a strong ranking suffices for excellent performance. The cost, however, is
$O(n^3)$---a full order of magnitude above classical additive decoders.

\paragraph{Standard GNN collapses under incremental decoding.}
Switching to the \change{PR} decoder exposes a sharp failure for SGNN under both
training regimes. Under objective training, incremental recovery falls to $0.354$ on
Easy and $0.285$ on Medium. This is the empirical signature of trajectory divergence
(Theorem~\ref{thm:propagation}): without fresh scores, small per-step errors
accumulate and the incremental path diverges from the rerun path.

\paragraph{Self-predictor training makes SGNN incrementally worse.}
The most informative result in Table~\ref{tab:main} is the further collapse of SGNN
under self-predictor training: incremental performance drops to $0.116$ on Easy and
$0.195$ on Medium---\emph{worse} than objective training alone ($0.354$ and $0.285$
respectively). This is not a pathology but a controlled diagnostic. The self-predictor
loss asks SGNN to predict the teacher's removal sequence from incrementally updated
scores. Without the gradient channel, no internal representation is available that
supports such updates: the imitation objective is unachievable given the architecture.
Training on an impossible signal does not help; it actively disrupts the useful static
ranking that objective training alone would have produced. This result makes a sharp
point: the gradient channel is not merely a helpful addition---it is a
\emph{necessary condition} for self-predictor training to be effective. Training
regime and architectural signal are interdependent; one without the other is either
neutral or harmful.

\paragraph{MAGNN closes the gap with the gradient channel alone.}
Even without self-predictor training, the gradient channel substantially improves
incremental decoding. MAGNN + Objective reaches $0.928$ on Easy and $0.772$ on
Medium---a $2.6\times$ and $2.7\times$ improvement over SGNN + Objective under the
same decoder. This confirms that access to local linear marginals, exposed via the
objective gradient, is itself sufficient to dramatically improve update compatibility,
even when the model is not explicitly trained to imitate the rerun trajectory.

\paragraph{MAGNN + Self achieves perfect incremental recovery.}
Combining the gradient channel with self-predictor training yields the strongest
result: MAGNN + Self achieves approximation $1.000$ and exact recovery $1.000$ on
both Easy and Medium under the \change{PR} decoder---matching the $O(n^3)$ rerun
baseline at $O(n^2)$ cost. On Hard, it reaches $0.689$ approximation, the best
incremental result in our comparison. The gradient channel makes the imitation
objective achievable; the imitation objective then drives the model to fully exploit
that signal.

\paragraph{Trade-off under self-predictor training.}
Table~\ref{tab:main} shows that MAGNN + Self has lower rerun performance on Medium
($0.698$) than MAGNN + Obj ($1.000$). This is expected: the self-predictor loss
pulls the model's internal representations toward update-compatible scores rather
than toward the best possible static ranking, slightly trading rerun quality for
incremental quality. Since the \change{PR} decoder is the primary target---and
achieves perfect recovery---this trade-off is well worthwhile.

\paragraph{Hard regime.}
Hard instances ($k \in [20, 35]$) remain difficult for all models under both
decoders. The best result is MAGNN + Self at $0.689$ approximation incremental.
Under the planted clique hardness conjecture~\citep{manurangsi2021}, no
polynomial-time algorithm is expected to fully solve this regime, so the residual
difficulty reflects a fundamental computational barrier rather than a failure of the
approach.

\subsection{One-Pass Static Ranking}
\label{sec:onepass}

Table~\ref{tab:onepass} reports the quality of the static ranking produced by each
model before any decoding. These results measure whether the GNN has learned a useful
global vertex ordering from a single forward pass, independently of the decoder.

\begin{table}[t]
  \centering
  \caption{One-pass static ranking performance (no decoding applied).}
  \label{tab:onepass}
  \begin{tabular}{lccccccccc}
    \toprule
    & \multicolumn{3}{c}{Easy} & \multicolumn{3}{c}{Medium}
    & \multicolumn{3}{c}{Hard} \\
    \cmidrule(lr){2-4}\cmidrule(lr){5-7}\cmidrule(lr){8-10}
    Model & Approx & Overlap & Exact
    & Approx & Overlap & Exact
    & Approx & Overlap & Exact \\
    \midrule
    SGNN + Obj
    & 0.964 & 0.963 & 0.900 & 0.692 & 0.637 & 0.400 & 0.422 & 0.046 & 0.000 \\
    MAGNN + Obj
    & 1.000 & 1.000 & 1.000 & 0.518 & 0.455 & 0.200 & 0.485 & 0.151 & 0.100 \\
    SGNN + Self
    & 1.000 & 1.000 & 1.000 & 0.505 & 0.438 & 0.200 & 0.375 & 0.118 & 0.000 \\
    MAGNN + Self
    & 0.963 & 0.961 & 0.900 & 0.810 & 0.797 & 0.600 & 0.448 & 0.140 & 0.000 \\
    \bottomrule
  \end{tabular}
\end{table}

All models learn meaningful static rankings on Easy instances. On Medium and Hard,
one-pass performance is weaker across the board, as expected: these regimes require
adaptive residual behavior that a fixed global ranking cannot provide. Notably, SGNN
+ Self achieves the best one-pass result on Easy ($1.000$ exact), yet collapses most
severely under incremental decoding. This is the clearest illustration that static
ranking quality and incremental update compatibility are orthogonal properties.
MAGNN's advantage is not in static ranking but in maintaining useful scores
throughout the pruning process.

\subsection{Embedding Geometry}
\label{sec:pca}

Table~\ref{tab:pca} reports PCA diagnostics on the final node embeddings.

\begin{table}[t]
  \centering
  \caption{PCA diagnostics on final node embeddings. PC1 and PC2 denote variance
    explained by the first two principal components. $\rho$ = Spearman correlation
  between original graph degree and the PC1 projection.}
  \label{tab:pca}
  \begin{tabular}{lccccccccc}
    \toprule
    & \multicolumn{3}{c}{Easy} & \multicolumn{3}{c}{Medium}
    & \multicolumn{3}{c}{Hard} \\
    \cmidrule(lr){2-4}\cmidrule(lr){5-7}\cmidrule(lr){8-10}
    Model & PC1 & PC2 & $\rho$
    & PC1 & PC2 & $\rho$
    & PC1 & PC2 & $\rho$ \\
    \midrule
    SGNN + Obj
    & 1.000 & 0.000 & $ 0.600$
    & 1.000 & 0.000 & $ 0.000$
    & 1.000 & 0.000 & $-0.200$ \\
    MAGNN + Obj
    & 0.997 & 0.003 & $ 0.400$
    & 0.997 & 0.003 & $-1.000$
    & 0.997 & 0.003 & $-1.000$ \\
    SGNN + Self
    & 1.000 & 0.000 & $ 1.000$
    & 1.000 & 0.000 & $-0.400$
    & 1.000 & 0.000 & $-1.000$ \\
    MAGNN + Self
    & 0.990 & 0.010 & $-0.200$
    & 0.989 & 0.011 & $-1.000$
    & 0.989 & 0.011 & $-1.000$ \\
    \bottomrule
  \end{tabular}
\end{table}

PC1 explains $>99\%$ of variance in all variants, confirming a dominant scalar
ordering direction and replicating the finding of~\citet{shoham2026ltr} that GNNs
collapse onto a degree-like principal component regardless of architecture.

The Spearman correlation $\rho$ varies in sign across models and regimes. The sign of
a principal component is arbitrary, so $\rho = -1$ is as strong a degree alignment as
$\rho = +1$. The structurally meaningful difference is elsewhere: SGNN embeddings are
strictly one-dimensional (PC2 $= 0.000$), whereas MAGNN embeddings retain a small
but non-zero second dimension (PC2 $\approx 0.003$--$0.010$). This residual structure
carries the geometric signature of the gradient channel---the marginal information
that $U_\theta$ exploits during incremental decoding. The fact that SGNN collapses to
a pure one-dimensional embedding even under self-predictor training explains why no
updater can rescue it: there is simply no second-order information for $U_\theta$ to
work with.

% ============================================================
%  DISCUSSION
% ============================================================
\section{Discussion}
\label{sec:discussion}

\paragraph{The core finding.}
The gap between $O(n^3)$ and $O(n^2)$ decoding is not a matter of approximation
quality or architectural expressiveness---it is a matter of \emph{update
compatibility}. A GNN can learn a perfect ranking and still fail under incremental
decoding, because its nonlinear encoding cannot be corrected by a fixed additive
rule. The gradient channel resolves this by exposing the local linear marginals of
the objective directly to the model, giving the learned updater the information it
needs to approximate the adaptive rerun trajectory.

\paragraph{The gradient channel is necessary, not just sufficient.}
The SGNN + Self result makes a point stronger than we might have anticipated: without
the gradient channel, adding an incremental training signal is actively harmful. The
self-predictor objective is well-defined and differentiable, yet it destroys SGNN's
incremental performance. This means the failure of standard GNNs is not a training
issue that better objectives could fix---it is a representational issue that requires
an architectural solution. The gradient channel is that solution.

\paragraph{Self-predictor training as a general design principle.}
The self-predictor regime is a general recipe for any differentiable combinatorial
objective: define a rerun-based teacher, train a student to imitate it incrementally,
and expose the objective gradient to make the imitation achievable. The teacher
provides the right removal sequence; the gradient channel provides the right
incremental signal; the updater learns the correction. The three components are
cleanly separable, suggesting they can be transferred to other problems where the
same $O(n^3)$-vs-$O(n^2)$ tension arises.

\paragraph{Limitations and failure modes.}
The Hard regime remains difficult for all models. At $k \in [20, 35]$ for $n=1000$,
the planted signal is weak and many dense distractor subgraphs exist; even MAGNN +
Self reaches only $0.689$ approximation. Under the planted clique hardness
conjecture, no polynomial-time algorithm is expected to fully solve this regime, so
the residual difficulty is a fundamental computational barrier. More practically, the
self-predictor loss requires running the teacher (a full rerun decoder) during
training, itself $O(n^3)$ per instance. This cost is paid at training time, not test
time, but it limits the practical size of training instances.

\paragraph{Relation to OptGNN and algorithmic alignment.}
OptGNN~\citep{yau2023} aligns GNN message passing with SDP gradient descent to
obtain approximation guarantees for Max-CSPs. Our work is complementary: we align
the internal signal of the GNN with the local marginal structure of the learned
objective to enable efficient residual decoding. Both approaches expose the right
optimization geometry to the network to improve performance. The difference is that
OptGNN fixes the target algorithm and designs the architecture to implement it, while
we fix the signal type and let the model learn the update rule.

\paragraph{Generality.}
Any combinatorial problem with a differentiable relaxation---\textsc{Max-Cut},
\textsc{Max-Independent-Set}, constraint satisfaction---has an objective gradient
that exposes additive local marginals. Inserting this gradient as an internal signal
and training a residual updater should enable $O(n^2)$ incremental decoding wherever
the $O(n^3)$ rerun decoder currently succeeds. We view this as the primary direction
for future work.

\paragraph{The algorithmic self-awareness principle.}
Classical combinatorial algorithms are efficient because they are \emph{self-aware}:
they maintain explicit state and update it exactly after each decision. Neural
networks trained only to rank lose this self-awareness---they produce scores useful
at one snapshot but unable to track the evolving state. Making the model self-aware
by feeding it the gradient of its own objective recovers a form of this classical
efficiency without sacrificing the learning-based flexibility that makes neural
solvers attractive. We believe this principle---exposing local optimization geometry
as an internal signal---is a promising direction for making neural combinatorial
algorithms efficient at scale.

\bibliographystyle{plainnat}
\bibliography{references}

\appendix

\section{Gradient Rectification: Formal Statement and Proof}
\label{app:rectification}

\change{We prove that the objective gradient $g_v$, when provided as an input to the
  score updater $U_\theta$, supplies exactly the information needed to recover a
  linear/additive feature in pre-activation space --- even though $U_\theta$ itself
  contains nonlinear activations. The key insight is that the combination of $g_v$ and
  $p_v$ determines the pre-activation gradient, which is the additive quantity that
  changes locally under each deletion. $U_\theta$'s linear first layer can extract this
  combination exactly before any nonlinearity is applied; subsequent layers then operate
on an already-additive signal.}

\begin{theorem}[Gradient in $U_\theta$ enables exact additive updates]
  \label{thm:rectification}
  \change{Let $\mathcal{L}(p) = \sum_r \alpha_r p^\top M_r p$ be a quadratic objective
    and let $g_i = \partial\mathcal{L}/\partial p_i$ be the per-vertex gradient. The
    correct score update preserving the removal ordering after deleting $v^*$ is to track
    \begin{equation}
      \Delta g_i = -2\sum_r \alpha_r (M_r)_{iv^*}\, p_{v^*},
      \label{eq:delta_g_app}
    \end{equation}
    which is a \emph{linear} function of inputs $(g_{v^*}, A_{iv^*}, p_{v^*})$. Then:
    \begin{enumerate}
      \item \textbf{With $g_i$ in $U_\theta$:} $\Delta g_i$ is a linear function of
        $U_\theta$'s inputs. A linear map --- the first layer of $U_\theta$ --- can
        represent $\Delta g_i$ exactly, before any nonlinear activation of $U_\theta$ is
        applied. Therefore $U_\theta$ can learn to track the correct ordering exactly,
        with no approximation error in the update rule.
      \item \textbf{Without $g_i$ in $U_\theta$:} $U_\theta$ receives only $(s_i,
        s_{v^*}, A_{iv^*})$ where $s_{v^*} = \sigma(z_{v^*})$ is a nonlinear encoding.
        Recovering $\Delta g_i$ from $s_{v^*}$ requires inverting $\sigma$, which by
        Theorem~\ref{thm:nonlinear} cannot be achieved with a fixed additive rule. The
        update is therefore an approximation, and by Theorem~\ref{thm:propagation} errors
        compound through the pruning trajectory.
    \end{enumerate}
    The gradient channel in $U_\theta$ thus converts an impossible approximation problem
    into a representable linear problem, for any monotone differentiable activation
  $\sigma$ and any quadratic objective.}
\end{theorem}

\begin{proof}
  \change{
    \textbf{Part 1.} The update $\Delta g_i$ in Eq.~\eqref{eq:delta_g_app} is linear in
    $p_{v^*}$ and in $(M_r)_{iv^*}$, both available as inputs $(p_{v^*}, A_{iv^*})$ to
    $U_\theta$. The full gradient $g_i$ is also an input, and $\Delta g_i$ does not
    depend on $g_i$ itself --- it depends on $g_{v^*}$ and $A_{iv^*}$. The first layer
    of $U_\theta$ computes $W[s_i, s_{v^*}, A_{iv^*}, g_i, g_{v^*}]$, a linear map.
    This layer can represent $\Delta g_i = -2\sum_r \alpha_r (M_r)_{iv^*} p_{v^*}$ by
    setting appropriate weights on $g_{v^*}$ and $A_{iv^*}$, before any nonlinear
    activation is applied. Subsequent layers receive a signal that is already the correct
    additive correction, and can learn to pass it through by approximating the identity
    --- which is in the representable class of any MLP.

    \textbf{Part 2.} Without $g_i$, the available inputs are $(s_i, s_{v^*}, A_{iv^*})$.
    Since $s_{v^*} = \sigma(z_{v^*})$ is a nonlinear function of the underlying statistic
    $d_{v^*}$, and by Theorem~\ref{thm:affine} the correct update requires affine
    dependence on $d_{v^*}$, recovering $\Delta g_i$ from $s_{v^*}$ requires inverting
    $\sigma$. By Theorem~\ref{thm:nonlinear}, no fixed additive rule can achieve this
    for all inputs. The update error at each step is therefore nonzero, and by
  Theorem~\ref{thm:propagation} these errors compound.}
\end{proof}

\begin{remark}
  \change{Note that the gradient in the GNN backbone (MAGNN without GU) does not help
    $U_\theta$: the backbone's gradient is computed on the \emph{original} graph at $t=0$
    and goes stale after the first deletion. $U_\theta$ needs the \emph{current}
    gradient, recomputed at each deletion step --- which it can compute from $(g_i,
    g_{v^*}, A_{iv^*}, p_{v^*})$ via Eq.~\eqref{eq:delta_g_app}. This explains the
    empirical finding that SGNN\_GU outperforms MAGNN\_U: it is the gradient in $U_\theta$
  that matters, not in the backbone.}
\end{remark}

\end{document}
